%%%%%%%% ICML 2026 EXAMPLE LATEX SUBMISSION FILE %%%%%%%%%%%%%%%%%

\documentclass{article}

% Recommended, but optional, packages for figures and better typesetting:
\usepackage{microtype}
\usepackage{graphicx}
\usepackage{subcaption}
\usepackage{booktabs} % for professional tables

% hyperref makes hyperlinks in the resulting PDF.
\usepackage{hyperref}

% Attempt to make hyperref and algorithmic work together better:
\newcommand{\theHalgorithm}{\arabic{algorithm}}

% Use the following line for the initial blind version submitted for review:
\usepackage{icml2026}

\usepackage{amsmath}
\usepackage{amssymb}
\usepackage{mathtools}
\usepackage{amsthm}

% if you use cleveref..
\usepackage[capitalize,noabbrev]{cleveref}

%%%%%%%%%%%%%%%%%%%%%%%%%%%%%%%%
% THEOREMS
%%%%%%%%%%%%%%%%%%%%%%%%%%%%%%%%
\theoremstyle{plain}
\newtheorem{theorem}{Theorem}[section]
\newtheorem{proposition}[theorem]{Proposition}
\newtheorem{lemma}[theorem]{Lemma}
\newtheorem{corollary}[theorem]{Corollary}
\theoremstyle{definition}
\newtheorem{definition}[theorem]{Definition}
\newtheorem{assumption}[theorem]{Assumption}
\theoremstyle{remark}
\newtheorem{remark}[theorem]{Remark}

\usepackage[textsize=tiny]{todonotes}

% The \icmltitle you define below is probably too long as a header.
\icmltitlerunning{LFWA}


% Clean spacing overrides to fit within strict page budget
\makeatletter
\renewcommand{\section}{\@startsection{section}{1}{\z@}%
                       {-8pt plus -2pt minus -2pt}%
                       {4pt plus 2pt}%
                       {\normalfont\large\bfseries\raggedright}}
\renewcommand{\subsection}{\@startsection{subsection}{2}{\z@}%
                          {-6pt plus -2pt minus -1pt}%
                          {3pt plus 1pt}%
                          {\normalfont\normalsize\bfseries\raggedright}}
\makeatother

\setlength{\textfloatsep}{8pt plus 2pt minus 2pt}
\setlength{\dbltextfloatsep}{8pt plus 2pt minus 2pt}
\setlength{\intextsep}{8pt plus 2pt minus 2pt}
\setlength{\dblfloatsep}{8pt plus 2pt minus 2pt}

\begin{document}
\setlength{\abovedisplayskip}{4pt}
\setlength{\belowdisplayskip}{4pt}
\setlength{\abovedisplayshortskip}{4pt}
\setlength{\belowdisplayshortskip}{4pt}

\twocolumn[
  \icmltitle{LFWA: Layer-wise Fisher-Weighted Adaptation for \\ Robust Test-Time Model Merging}

  \icmlsetsymbol{equal}{*}

  \begin{icmlauthorlist}
    \icmlauthor{Anonymous Authors}{equal}
  \end{icmlauthorlist}

  \icmlaffiliation{anon}{Anonymous Institution, Anonymous City, Anonymous Country}

  \icmlkeywords{Machine Learning, Model Merging, Test-time Adaptation}

  \vskip 0.3in
]

\printAffiliationsAndNotice{}

\makeatletter
\gdef\@icmltitlerunning{LFWA: Layer-wise Fisher-Weighted Test-Time Model Merging}
\makeatother

\begin{abstract}
  Model merging has emerged as a highly cost-effective paradigm to integrate specialized capabilities of independently fine-tuned expert models into a single multi-task system without requiring joint training. To adapt to dynamic downstream test environments, recent methods employ online test-time adaptation (TTA) of merging coefficients via unsupervised entropy minimization. However, these methods use a uniform global learning rate across all layers, which ignores the highly heterogeneous sensitivity of neural network representations. In this paper, we propose \textbf{Layer-wise Fisher-Weighted Adaptation (LFWA)}, a novel framework that stabilizes and accelerates test-time model merging. LFWA utilizes pre-computed average parameter-level diagonal Fisher Information as a sensitivity prior to scale layer-wise learning rates dynamically, damping updates in highly sensitive layers (such as early convolutional blocks and classification heads) and accelerating them in robust representational layers. Evaluated on both alternating and block-sequential non-stationary streams of CIFAR-10 and SVHN tasks using a pre-trained ResNet-18, LFWA achieves outstanding performance and shows exceptional stability, outperforming static merging and standard uniform-rate TTA across both Adam and SGD optimizers. Furthermore, LFWA substantially reduces sensitivity to the global learning rate, providing a highly robust and practical optimization foundation for online adaptive model merging.
\end{abstract}

\section{Introduction}
\label{sec:intro}

The pre-training and fine-tuning paradigm has driven deep learning to massive success across diverse applications \cite{devlin2018bert,vaswani2017attention,he2016deep,dosovitskiy2020image,aslam2024continual,baysal2025overcoming,tseng2025incremental}. However, fine-tuning large base models independently on multiple downstream tasks yields many separate experts, scaling parameter storage and hosting costs linearly with the number of tasks. Joint multi-task training addresses this but is often computationally prohibitive, requires simultaneous access to all private datasets, and introduces severe optimization difficulties such as gradient conflicts and task-balancing issues \cite{yadav2023ties,mao2025multi-task,wu2025attention-based,cao2024multi-task,huang2025optimization,boiarov2021simultaneous,he2024robust,fan2023maxgnr,choi2023task-aware}.

To circumvent these limitations, model merging (or deep model fusion) has emerged as an attractive, training-free alternative \cite{wortsman2022model,ilharco2022editing,li2023deep,tang2024fusionbench,saadati2024dimat,du2024parameter,liu2025sens-merging,qiu2026parameter,zeng2025robustmerge,zeng2025parameter}. Model merging directly combines independently fine-tuned weights at the parameter level. Traditional methods such as Task Arithmetic \cite{ilharco2022editing} and TIES-Merging \cite{yadav2023ties} perform simple linear operations in Euclidean space to merge task-specific vectors. While highly efficient, these Euclidean averaging techniques suffer from severe parameter interference, where conflicting updates from different experts cancel each other out, degrading the merged model's performance on downstream tasks. Recent works have extended parameter merging to large language models, multimodal models, and robot policies \cite{jang2023personalized,yadav2025robust,li2025multi-modality,hui2024upcycling,ju2024mitigating}.

To address parameter interference, test-time adaptive merging methods, such as AdaMerging \cite{yang2023adamerging}, optimize merging coefficients on unlabeled test data via prediction entropy minimization. More recently, SyMerge \cite{jung2025symerge} redefined the goal from mitigating interference to fostering active cross-task synergy by jointly adapting classification heads and layer-wise merging coefficients. However, unsupervised test-time adaptation (TTA) of merging coefficients is highly susceptible to overfitting, parameter drift, and representation collapse under severe out-of-distribution (OOD) shifts and spatial corruptions \cite{huang2022extrapolative,zhao2025threshold,lei2025two-level,jia2024tinytta,yeo2023rapid,bhardwaj2022benchmarking,zhang2024unsupervised,chen2022contrastive,kim2023when,zhang2025scap,gui2024active,yin2025test-time}.

A major driver of this instability is that existing TTA model merging frameworks employ a uniform global learning rate across all layers. This ignores a fundamental property of deep neural networks: different layers exhibit highly heterogeneous representational and structural sensitivity. For example, early feature-extracting layers and final classification heads are highly sensitive to parameter changes, whereas intermediate representational layers are relatively robust. Updating merging coefficients uniformly across all layers leads to a destructive trade-off: a learning rate large enough to adapt robust layers causes catastrophic representation collapse in sensitive layers, whereas a learning rate small enough to protect sensitive layers stalls adaptation in others.

In this work, we propose \textbf{Layer-wise Fisher-Weighted Adaptation (LFWA)}, a novel, mathematically principled optimization framework that stabilizes and accelerates online test-time model merging. LFWA utilizes pre-computed average parameter-level diagonal Fisher Information as a sensitivity prior to scale layer-wise learning rates dynamically. Specifically, the learning rate for each layer's merging coefficient is scaled inversely proportional to that layer's average parameter Fisher sensitivity. This ensures that sensitive blocks undergo small, controlled adaptations to prevent representation collapse, while robust layers can adapt rapidly to dynamic test distributions.

Our contributions are summarized as follows:
\begin{itemize}\setlength{\itemsep}{1pt}\setlength{\parskip}{0pt}\setlength{\parsep}{0pt}
    \item We propose LFWA, a parameter-efficient, zero-inference-overhead optimization framework that incorporates layer-specific Fisher Information priors to stabilize test-time model merging.
    \item We demonstrate that LFWA resolves the fundamental trade-off of uniform-rate TTA, enabling safe, high-velocity adaptation without representation collapse.
    \item We provide a mathematical formulation connecting LFWA to Natural Gradient Descent (NGD), showing that LFWA acts as an elegant block-diagonal preconditioning operator.
    \item Through extensive sweeps on alternating and block-sequential streams of CIFAR-10 and SVHN using a pre-trained ResNet-18, we show that LFWA consistently outperforms static merging and standard TTA across both Adam and SGD optimizers, while dramatically reducing sensitivity to the global learning rate.
\end{itemize}

\section{Related Work}
\label{sec:related}

\subsection{Model Merging and Deep Model Fusion}
Model merging aims to fuse multiple neural networks fine-tuned from a shared pre-trained checkpoint into a single unified model. Early work focused on simple parameter averaging (Model Soups) \cite{wortsman2022model} and Task Arithmetic \cite{ilharco2022editing}, which computes task-specific vectors and adds them to the base model. To resolve parameter conflicts, TIES-Merging \cite{yadav2023ties} prunes redundant parameters and resolves sign disagreements, while DARE \cite{yu2024dare} uses random dropouts to sparsify task vectors. Recent surveys have categorized these methods under deep model fusion \cite{li2023deep,tang2024fusionbench}. Advanced algorithms such as DIMAT \cite{saadati2024dimat}, Parameter Competition Balancing \cite{du2024parameter}, and sensitivity-guided balancing \cite{liu2025sens-merging} have improved the trade-offs. Furthermore, parameter merging has expanded into domain-specific applications like robots, multi-modality expansion, MoE upcycling, and LLM alignment \cite{jang2023personalized,yadav2025robust,li2025multi-modality,hui2024upcycling,ju2024mitigating,zeng2025robustmerge,zeng2025parameter,qiu2026parameter}. However, these methods are static and cannot adapt to dynamic distribution shifts in the test stream.

\subsection{Test-Time Adaptation of Deep Networks}
Test-time adaptation (TTA) adjusts a pre-trained model to a shifting target domain using unlabeled test streams. Unsupervised TTA algorithms like Tent \cite{chen2022contrastive} use entropy minimization, while others leverage contrastive learning, time-series anomaly detection, and early-exit ensembles \cite{chen2022contrastive,kim2023when,jia2024tinytta}. TTA has been extended to spiking networks, multi-modal learning, gaze estimation, and active feedback systems \cite{zhao2025threshold,lei2025two-level,yin2025test-time,yeo2023rapid,gui2024active,huang2022extrapolative,bhardwaj2022benchmarking,zhang2024unsupervised,zhang2025scap}. In the context of merged models, AdaMerging \cite{yang2023adamerging} optimizes layer-wise weights on-the-fly using prediction entropy minimization on target batches. SyMerge \cite{jung2025symerge} and SATA-SBF / SATA-RGP \cite{submission1} optimize both coefficients and heads using expert-guided self-labeling/teacher guidance, which requires keeping all unmerged expert models (teachers) in memory, resulting in massive VRAM and compute overhead. EWC-TTA \cite{submission7} guides dynamic model merging under task-switching using pre-computed diagonal Fisher priors as a quadratic penalty on head weights. Recent teacher-free methods like S2C-Merge \cite{submission5} freeze classification heads to prevent decision-boundary collapse. In contrast, our proposed LFWA represents a complementary, foundational improvement to the underlying optimization process, using Fisher sensitivity priors to safely adapt all layers without collapse.

\subsection{Natural Gradient and Fisher Preconditioning}
Natural Gradient Descent (NGD) \cite{amari1998natural} scales parameter updates by the inverse of the Fisher Information Matrix (FIM) to follow the steepest descent direction in probability space. While NGD offers superior optimization properties, computing and inverting the full FIM is computationally intractable. Consequently, researchers have proposed diagonal and block-diagonal approximations, Kronecker-factored curvature, components-wise NGD, randomized linear algebra, and dual NGD formulations to scale NGD to deep networks and physics-informed models \cite{rudner2019natural,liu2024reconstructing,bioli2025accelerating,bao2025projected,abdulkadirov2023solving,sang2022component-wise,jnini2025dual,jnini2024gauss-newton}. Our proposed LFWA establishes a theoretical bridge between TTA model merging and NGD, leveraging pre-computed average parameter-level diagonal Fisher sensitivities as a lightweight block-diagonal preconditioning operator on the merging coefficient space.

\subsection{Elastic Weight Consolidation and Sharpness-Aware Minimization}
Elastic Weight Consolidation (EWC) uses diagonal Fisher Information priors as a quadratic penalty to preserve important parameters during sequential task learning, mitigating catastrophic forgetting \cite{he2025robust,xue2025enhancing,jhajj2025elastic,aslam2024continual,li2020few-shot,baysal2025overcoming,tseng2025incremental,sliogeris2025elastic,liu2025uncertainty-aware,wang2025development}. Concurrently, Sharpness-Aware Minimization (SAM) improves generalization by optimizing for both loss value and loss sharpness, restricting parameters from settling in narrow, sensitive minima \cite{foret2020sharpness-aware,andriushchenko2022towards,li2024friendly,ilbert2024samformer,fan2025towards,luo2025explicit,li2025focal-sam,oikonomou2025sharpness-aware,qu2022generalized}. While EWC and SAM focus on training parameters, LFWA uniquely applies Fisher Information to modulate the optimization dynamics of *merging coefficients* during online test-time adaptation, avoiding representation collapse without adding training penalties or extra optimization steps.

\section{Methodology}
\label{sec:methodology}

We now describe our proposed framework, \textbf{Layer-wise Fisher-Weighted Adaptation (LFWA)}, for online test-time model merging.

\subsection{Preliminaries and Problem Formulation}
Let $f_{\theta_{\text{pre}}}$ represent a base model encoder parameterized by weights $\theta_{\text{pre}} \in \mathbb{R}^P$. We are given a set of $K$ expert models, where each expert $k$ shares the same base architecture but has been independently fine-tuned on a task-specific dataset, yielding encoder parameters $\theta_k \in \mathbb{R}^P$ and a task-specific classification head $h_{\phi_k}$ parameterized by $\phi_k$. Following Task Arithmetic \cite{ilharco2022editing}, we define the task vector for expert $k$ as $\tau_k = \theta_k - \theta_{\text{pre}}$.

Model merging combines these experts by constructing a virtual merged encoder parameterized by $\Theta(\lambda)$, where:
\begin{equation}
\Theta(\lambda) = \theta_{\text{pre}} + \sum_{k=1}^K \lambda_{k} \tau_k
\end{equation}
and $\lambda_k \in \mathbb{R}$ represents the merging coefficient for task $k$. To allow for fine-grained representational flexibility, we optimize \emph{layer-wise} (or tensor-wise) merging coefficients, where each named parameter tensor $w$ has task-specific coefficients $\lambda_{k, w}$, yielding merged tensor parameters:
\begin{equation}
w_{\text{merged}}(\lambda) = w_{\text{base}} + \sum_{k=1}^K \lambda_{k, w} (w_{\text{expert}_k} - w_{\text{base}})
\end{equation}

During test-time adaptation (TTA), the merged model is deployed on an unlabeled, non-stationary test stream. For each incoming batch $X_t$ from a dynamic task context $k$, we evaluate soft-predictions using the merged encoder $\Theta(\lambda)$ and the active head $h_{\phi_k}$:
\begin{equation}
P_{\text{merged}} = \text{Softmax}\left(h_{\phi_k}\left(f_{\Theta(\lambda)}(X_t)\right)\right)
\end{equation}
The merging coefficients are dynamically optimized on-the-fly to minimize prediction entropy:
\begin{equation}
\label{eq:entropy_loss}
\begin{split}
\mathcal{L}_{\text{ent}}(X_t; \lambda) = &- \frac{1}{|X_t|} \sum_{x \in X_t} \sum_{c} P_{\text{merged}}(x)_c \\
&\cdot \log\left(P_{\text{merged}}(x)_c + \epsilon_{\text{ent}}\right)
\end{split}
\end{equation}
using a gradient descent update:
\begin{equation}
\lambda_{k, w} \leftarrow \lambda_{k, w} - \eta_{w} \nabla_{\lambda_{k, w}} \mathcal{L}_{\text{ent}}
\end{equation}
where $\eta_w$ is the learning rate assigned to tensor $w$'s coefficients.

\subsection{Layer-wise Parameter-level Fisher Sensitivity Estimation}
Standard TTA model merging employs a homogeneous, uniform learning rate across all tensors ($\eta_w = \eta$). However, different parameter tensors in deep networks exhibit vastly different sensitivities. To quantify this sensitivity, we pre-compute the diagonal empirical Fisher Information Matrix (FIM) on a tiny calibration subset of $N_{\text{cal}}$ samples for each expert model $k$ before deployment.

For a parameter $p$ in tensor $w$, the diagonal empirical Fisher Information is defined as:
\begin{equation}
F_{p} = \frac{1}{N_{\text{cal}}} \sum_{i=1}^{N_{\text{cal}}} \left(\frac{\partial \mathcal{L}_{\text{CE}}(x_i)}{\partial \theta_p}\right)^2
\end{equation}
where $\mathcal{L}_{\text{CE}}$ is the task-specific cross-entropy loss. To obtain a robust, low-overhead scalar sensitivity prior for each named parameter tensor $w$, we average the diagonal Fisher values over all individual parameters $p$ within that tensor:
\begin{equation}
\bar{F}_{k, w} = \frac{1}{|w|} \sum_{p \in w} F_{p}
\end{equation}
We then compute the joint layer-wise Fisher sensitivity $\bar{F}_w$ across all experts:
\begin{equation}
\bar{F}_w = \frac{1}{K} \sum_{k=1}^K \bar{F}_{k, w}
\end{equation}
This joint layer-wise sensitivity requires storing only a single scalar per tensor (yielding less than 30 floats for a ResNet-18 model), adding virtually zero storage or memory overhead.

\subsection{Fisher-Weighted Adaptive Learning Rates}
Using the joint layer-wise Fisher sensitivity prior $\bar{F}_w$, we propose to define heterogeneous, sensitivity-aware learning rates for each tensor's merging coefficients:
\begin{equation}
\eta_{w} = \frac{\eta}{\left(\bar{F}_w + \epsilon_{\text{scale}}\right)^\alpha}
\end{equation}
where $\eta$ is the global learning rate, $\epsilon_{\text{scale}}$ is a small constant to prevent division by zero, and $\alpha \ge 0$ is the sensitivity damping hyperparameter. 

The scaling power $\alpha$ acts as a continuous dial of optimization heterogeneity:
\begin{itemize}\setlength{\itemsep}{1pt}\setlength{\parskip}{0pt}\setlength{\parsep}{0pt}
    \item When $\alpha = 0$, we recover standard uniform-rate TTA (such as AdaMerging), where all layers adapt at rate $\eta$.
    \item When $\alpha > 0$, the learning rates are adaptively damped in highly sensitive blocks (large $\bar{F}_w$, such as early feature extractors or classification-adjacent layers) and accelerated in robust, deep representational blocks (small $\bar{F}_w$).
\end{itemize}
By damping updates in sensitive parameters, LFWA prevents catastrophic representational drift and decision-boundary collapse. Simultaneously, by maintaining higher learning rates in less sensitive layers, LFWA encourages fast and efficient adaptation to shifting distributions, resolving the fundamental speed-stability trade-off of uniform-rate TTA.

\subsection{Theoretical Connection to Natural Gradient Descent}
We can formally justify LFWA by analyzing its connection to Natural Gradient Descent (NGD) \cite{amari1998natural} and demonstrating how parameter-level diagonal Fisher Information induces an elegant block-diagonal preconditioning operator on the space of merging coefficients.

In standard deep network optimization, NGD updates parameters along the steepest descent direction in the probability distribution space rather than the Euclidean parameter space. The NGD update rule is defined as:
\begin{equation}
\theta^{(t+1)} \leftarrow \theta^{(t)} - \eta F_{\theta}^{-1} \nabla_\theta \mathcal{L}(\theta)
\end{equation}
where $F_\theta \in \mathbb{R}^{P \times P}$ is the Fisher Information Matrix (FIM). Under linear model merging, the virtual merged parameters are a linear function of the merging coefficients $\lambda \in \mathbb{R}^{K \times L}$. Specifically, for a parameter $p$ in tensor $w$, the virtual parameter is:
\begin{equation}
\theta_p(\lambda) = w_{\text{base}, p} + \sum_{j=1}^K \lambda_{j, w} \tau_{j, p}
\end{equation}
where $\tau_{j, p} = w_{\text{expert}_j, p} - w_{\text{base}, p}$ is the task vector component of expert $j$. Using the chain rule, the gradient of the loss $\mathcal{L}$ with respect to the merging coefficient $\lambda_{k, w}$ for expert $k$ at layer $w$ is:
\begin{equation}
\frac{\partial \mathcal{L}}{\partial \lambda_{k, w}} = \sum_{p \in w} \frac{\partial \mathcal{L}}{\partial \theta_p} \frac{\partial \theta_p}{\partial \lambda_{k, w}} = \sum_{p \in w} \frac{\partial \mathcal{L}}{\partial \theta_p} \tau_{k, p}
\end{equation}
This shows that the gradient in the coefficient space is the projection of the parameter-space gradients along the task vector direction.

To find the induced Fisher Information Matrix $F_{\lambda}$ in the coefficient space, we compute the expectation of the outer product of the gradients with respect to the coefficients. For a specific layer $w$ and experts $k$ and $m$, the corresponding entry in the coefficient-space FIM is:
\begin{equation}
\begin{split}
(F_{\lambda})_{k, m} &= \mathbb{E} \left[ \left(\frac{\partial \mathcal{L}}{\partial \lambda_{k, w}}\right) \left(\frac{\partial \mathcal{L}}{\partial \lambda_{m, w}}\right) \right] \\
&= \mathbb{E} \left[ \left( \sum_{p \in w} \frac{\partial \mathcal{L}}{\partial \theta_p} \tau_{k, p} \right) \left( \sum_{q \in w} \frac{\partial \mathcal{L}}{\partial \theta_q} \tau_{m, q} \right) \right]
\end{split}
\end{equation}
Expanding this product, we obtain:
\begin{equation}
(F_{\lambda})_{k, m} = \sum_{p \in w} \sum_{q \in w} \mathbb{E} \left[ \frac{\partial \mathcal{L}}{\partial \theta_p} \frac{\partial \mathcal{L}}{\partial \theta_q} \right] \tau_{k, p} \tau_{m, q}
\end{equation}

By invoking a standard block-diagonal approximation of the parameter-space FIM where cross-parameter sensitivities within a layer are assumed to be approximately uncorrelated (i.e., $\mathbb{E} \left[ \frac{\partial \mathcal{L}}{\partial \theta_p} \frac{\partial \mathcal{L}}{\partial \theta_q} \right] \approx 0$ for $p \neq q$), the double summation simplifies to a single sum over the diagonal parameter-level Fisher Information $F_p = \mathbb{E} \left[ \left(\frac{\partial \mathcal{L}}{\partial \theta_p}\right)^2 \right]$:
\begin{equation}
(F_{\lambda})_{k, m} \approx \sum_{p \in w} F_p \tau_{k, p} \tau_{m, p}
\end{equation}

In the case where we evaluate the diagonal entries ($k = m$) of the coefficient-space FIM for task $k$, we have:
\begin{equation}
(F_{\lambda})_{k, k} \approx \sum_{p \in w} F_p \tau_{k, p}^2
\end{equation}
We can express the task vector component squares as $\tau_{k, p}^2 = |w| \cdot \text{Var}_p(\tau_{k, p}) + \mathbb{E}_p[\tau_{k, p}]^2$. Under the mild assumption that the average magnitude of task vector components is relatively homogeneous across the layer, we can approximate the summation by replacing the individual $\tau_{k, p}^2$ components with their average squared magnitude $\bar{\tau}_{k, w}^2 = \frac{1}{|w|} \sum_{p \in w} \tau_{k, p}^2$. This yields:
\begin{equation}
(F_{\lambda})_{k, k} \approx \bar{\tau}_{k, w}^2 \sum_{p \in w} F_p = |w| \bar{\tau}_{k, w}^2 \bar{F}_{k, w}
\end{equation}
where $\bar{F}_{k, w} = \frac{1}{|w|} \sum_{p \in w} F_p$ is precisely the average parameter-level diagonal Fisher Information of layer $w$ for expert $k$.

Taking the joint sensitivity across all $K$ expert models, we define the layer-wise preconditioning scaling factor $\bar{F}_w = \frac{1}{K} \sum_{k=1}^K \bar{F}_{k, w}$.
By preconditioning the gradient updates of the merging coefficients $\lambda_{k, w}$ with the inverse of the joint layer-wise Fisher sensitivity $\bar{F}_w$, we obtain:
\begin{equation}
\lambda_{k, w}^{(t+1)} \leftarrow \lambda_{k, w}^{(t)} - \frac{\eta}{\left(\bar{F}_w + \epsilon_{\text{scale}}\right)^\alpha} \nabla_{\lambda_{k, w}} \mathcal{L}
\end{equation}

This derivation provides a rigorous mathematical bridge: scaling the learning rate of the layer-wise merging coefficients inversely proportional to the average parameter-level Fisher Information serves as a direct, first-order approximation to Natural Gradient Descent on the coefficient space. Setting $\alpha = 1.0$ corresponds to full diagonal Fisher preconditioning, while $\alpha \in (0, 1)$ smoothly interpolates between standard gradient descent in Euclidean space and natural gradient updates in the probability distribution space. This formalizes why LFWA dramatically stabilizes the adaptation process: it scales the learning step sizes to match the information-theoretic geometry of the model parameters, preventing representation collapse in highly sensitive layers while accelerating adaptation in robust ones.

\section{Experimental Setup}
\label{sec:experiments}

To evaluate our proposed Layer-wise Fisher-Weighted Adaptation (LFWA) framework, we construct a non-stationary test-time adaptive model merging benchmark and compare LFWA against standard uniform-rate test-time adaptation and static model merging.

\subsection{Benchmark and Datasets}
Following test-time model merging literature \cite{yang2023adamerging,jung2025symerge}, we evaluate on a multi-task learning challenge consisting of two distinct image classification tasks: CIFAR-10 \cite{krizhevsky2009learning} and SVHN \cite{netzer2011reading}.
Each dataset represents a unique image domain (natural objects vs. house number digits) but shares the same output dimensionality of 10 classes.
We fine-tune separate experts starting from a shared pre-trained ResNet-18 initialization:
\begin{itemize}\setlength{\itemsep}{1pt}\setlength{\parskip}{0pt}\setlength{\parsep}{0pt}
    \item \textbf{Expert 1 (CIFAR-10)}: Fine-tuned on a subset of 5,000 CIFAR-10 training images for 3 epochs, achieving a standalone accuracy of 98.56\%.
    \item \textbf{Expert 2 (SVHN)}: Fine-tuned on a subset of 5,000 SVHN training images for 3 epochs, achieving a standalone accuracy of 93.60\%.
\end{itemize}

\subsection{Non-Stationary Test Streams}
To evaluate the impact of task-switching dynamics on test-time model merging, we construct two different test streams, each containing 32 batches (batch size 64) for a total of 2,048 test images:
\begin{enumerate}\setlength{\itemsep}{1pt}\setlength{\parskip}{0pt}\setlength{\parsep}{0pt}
    \item \textbf{Alternating Stream (High Frequency)}: Interleaves evaluation batches such that the model receives a stream of alternating task batches: Batch 1 (CIFAR-10), Batch 2 (SVHN), Batch 3 (CIFAR-10)... This simulates high-frequency, sudden domain-switching environments.
    \item \textbf{Block-Sequential Stream (Low Frequency)}: Processes all 16 batches of CIFAR-10 first, followed by all 16 batches of SVHN. This represents low-frequency, gradual task transition environments.
\end{enumerate}
Task classification heads are assigned based on the task context of the incoming batch, and only the 24 layer-wise merging coefficients in the encoder are updated via test-time adaptation.

\subsection{Baselines}
We compare our proposed method against two standard baselines:
\begin{enumerate}\setlength{\itemsep}{1pt}\setlength{\parskip}{0pt}\setlength{\parsep}{0pt}
    \item \textbf{Static Merging (Task Arithmetic Baseline)}: The merging coefficients are held fixed at a uniform initialization of $\lambda_{1, w} = 0.5$ and $\lambda_{2, w} = 0.5$ across all layers throughout the stream, representing a static, non-adaptive merged model.
    \item \textbf{Standard TTA (AdaMerging Baseline)}: The layer-wise merging coefficients are adapted on-the-fly using a uniform, homogeneous learning rate ($\eta_w = \eta$, corresponding to $\alpha = 0$ in our framework) via prediction entropy minimization on each test batch.
\end{enumerate}

\subsection{Optimization and Hyperparameters}
We evaluate the generalization of LFWA across two different optimization algorithms: **Adam** and **SGD with Momentum (0.9)**.
All methods perform a single gradient descent step on each incoming test batch before making predictions. The merging coefficients are clamped to $[0.0, 1.0]$ after each step to keep them within a physically grounded range. We evaluate a comprehensive grid of global learning rates $\eta \in \{0.001, 0.01, 0.1, 1.0\}$ and sensitivity powers $\alpha \in \{0.0, 0.2, 0.5, 1.0\}$. We set $\epsilon_{\text{scale}} = 10^{-8}$ for stable learning rate scaling.

\section{Empirical Results and Discussion}
\label{sec:results}

Our main experimental results are summarized in Table~\ref{tab:main_results}.

\begin{table*}[t]
\caption{Multi-task test-time adaptation performance (Accuracy \%) under both **Alternating (High Frequency)** and **Block-Sequential (Low Frequency)** CIFAR-10 / SVHN streams using ResNet-18. We compare Static Merging, Standard TTA, and our proposed LFWA across Adam and SGD optimizers. Best results for each optimizer-stream combination are highlighted in \textbf{bold}.}
\label{tab:main_results}
\vskip 0.15in
\begin{center}
\begin{small}
\resizebox{\textwidth}{!}{%
\begin{tabular}{lccccc}
\toprule
& & \multicolumn{2}{c}{\textbf{Alternating Stream (High Frequency) (\%)}} & \multicolumn{2}{c}{\textbf{Block-Sequential Stream (Low Frequency) (\%)}} \\
\cmidrule(r){3-4} \cmidrule(l){5-6}
Method & Learning Rate $\eta$ & Adam & SGD & Adam & SGD \\
\midrule
Static Merging Baseline & -- & 74.38 & 74.38 & 74.38 & 74.38 \\
\midrule
Standard TTA ($\alpha = 0.0$) & 0.001 & 74.73 & 75.29 & 74.82 & 74.79 \\
Standard TTA ($\alpha = 0.0$) & 0.01  & 78.06 & 77.50 & 77.56 & 76.47 \\
Standard TTA ($\alpha = 0.0$) & 0.1   & 78.82 & 78.48 & 81.26 & 77.96 \\
Standard TTA ($\alpha = 0.0$) & 1.0   & 78.66 & 79.02 & 80.95 & 81.34 \\
\midrule
LFWA ($\alpha = 0.2$)         & 0.001 & 77.76 & 77.67 & 77.58 & 76.43 \\
LFWA ($\alpha = 0.2$)         & 0.01  & 78.77 & 77.66 & 81.17 & 77.91 \\
LFWA ($\alpha = 0.2$)         & 0.1   & 79.17 & \textbf{79.40} & \textbf{81.55} & 81.83 \\
LFWA ($\alpha = 0.2$)         & 1.0   & 78.57 & 79.02 & 79.79 & \textbf{81.86} \\
\midrule
LFWA ($\alpha = 0.5$)         & 0.001 & \textbf{79.67} & 78.64 & 81.05 & 80.12 \\
LFWA ($\alpha = 0.5$)         & 0.01  & 79.14 & \textbf{79.46} & 80.74 & 81.13 \\
LFWA ($\alpha = 0.5$)         & 0.1   & 78.12 & 78.64 & 79.83 & 80.83 \\
LFWA ($\alpha = 0.5$)         & 1.0   & 76.48 & 75.72 & 79.57 & 78.47 \\
\midrule
LFWA ($\alpha = 1.0$)         & 0.001 & 77.72 & 79.01 & 80.21 & 81.19 \\
LFWA ($\alpha = 1.0$)         & 0.01  & 77.73 & 77.82 & 78.49 & 80.43 \\
LFWA ($\alpha = 1.0$)         & 0.1   & 76.03 & 76.76 & 79.78 & 79.63 \\
LFWA ($\alpha = 1.0$)         & 1.0   & 76.81 & 77.63 & 77.96 & 79.70 \\
\bottomrule
\end{tabular}%
}
\end{small}
\end{center}
\end{table*}

\subsection{Quantitative Performance Comparison}
As seen in the results, the non-adaptive Static Merging baseline yields an average multi-task accuracy of $74.38\%$. Online TTA successfully improves on this baseline: under standard uniform-rate TTA ($\alpha=0.0$), the Adam optimizer achieves a peak accuracy of $78.82\%$ in the Alternating stream (at $\eta=0.1$) and $81.26\%$ in the Sequential stream (at $\eta=0.1$). 

Our proposed LFWA framework consistently outperforms standard TTA across both streams and both optimizers. In the Alternating stream, Adam LFWA with $\alpha=0.5$ achieves a peak accuracy of \textbf{79.67\%} at $\eta=0.001$. In the Block-Sequential stream, Adam LFWA with $\alpha=0.2$ achieves a peak accuracy of \textbf{81.55\%} at $\eta=0.1$. Similarly, for the SGD optimizer, LFWA reaches a peak of \textbf{79.46\%} in the Alternating stream ($\alpha=0.5$, $\eta=0.01$) and \textbf{81.86\%} in the Sequential stream ($\alpha=0.2$, $\eta=1.0$), demonstrating the robust generalization of our Fisher preconditioning across optimization algorithms.

\subsection{Impact of Task-Switching Frequency}
Comparing the results in Table~\ref{tab:main_results} reveals an important empirical insight: test-time model adaptation yields significantly higher overall performance when the test stream transitions gradually (Sequential stream, up to $81.86\%$) compared to when it switches rapidly (Alternating stream, up to $79.67\%$). 

Under the Sequential stream, the model processes large homogenous blocks of a single task, allowing the merging coefficients to settle into optimal regimes for that task. Under the Alternating stream, the task changes on every single batch, causing the merging coefficients to oscillate rapidly. In this high-frequency oscillation regime, LFWA's sensitivity damping plays an even more critical role: by restricting updates in sensitive parameters, it prevents the coefficients of sensitive layers from "thrashing", which stabilizes the shared feature extraction pathway and guarantees robust representation quality.

\subsection{Dramatically Accelerated Adaptation at Low Learning Rates}
The most prominent benefit of LFWA is visible at low learning rates ($\eta = 0.001$). Under small global learning rates, standard uniform-rate TTA adapts too slowly and barely improves over the static baseline ($74.73\%$ vs. $74.38\%$ for Adam Alternating, and $74.79\%$ vs $74.38\%$ for SGD Sequential). 
By incorporating Fisher sensitivity priors, LFWA dynamically scales up the learning rates of representational layers that can safely absorb larger updates. For example:
\begin{itemize}\setlength{\itemsep}{1pt}\setlength{\parskip}{0pt}\setlength{\parsep}{0pt}
    \item In the Alternating stream, Adam LFWA ($\alpha=0.5$) boosts accuracy from $74.73\%$ to \textbf{79.67\%} (a massive \textbf{+4.94\%} absolute gain).
    \item In the Sequential stream, Adam LFWA ($\alpha=0.5$) boosts accuracy from $74.82\%$ to \textbf{81.05\%} (a massive \textbf{+6.23\%} absolute gain).
    \item Similarly, SGD LFWA ($\alpha=1.0$) on the Sequential stream boosts accuracy from $74.79\%$ to \textbf{81.19\%} (a massive \textbf{+6.40\%} absolute gain).
\end{itemize}
This strongly validates our hypothesis: LFWA successfully decouples layer-wise adaptation speeds, enabling rapid representational shift without risk of parameter collapse.

\subsection{Mitigating Sensitivity to the Global Learning Rate}
Unsupervised TTA is notoriously sensitive to hyperparameters, where a high learning rate causes rapid representation collapse, and a small rate leads to frozen weights. Standard uniform TTA is highly dependent on precise tuning of the global learning rate (for instance, ranging from $74.79\%$ to $81.34\%$ for SGD Sequential). 
In contrast, LFWA exhibits exceptional self-calibrating behavior and hyperparameter robustness. With $\alpha=0.2$, Adam maintains high, stable performance between $77.76\%$ and $79.17\%$ in the Alternating stream, and between $77.58\%$ and $81.55\%$ in the Sequential stream across learning rates spanning three orders of magnitude. This makes LFWA an exceptionally practical framework for real-world deployments where labeled validation sets are unavailable for hyperparameter sweeps.


\section{Limitations and Future Work}
\label{sec:limitations}
While LFWA shows remarkable performance, it has a few limitations:
First, estimating the joint Fisher Information requires a small calibration dataset from the expert training distribution. Although this calibration set is tiny (500 images) and does not need to be stored after pre-computation, it may not be available if training data is completely private. However, as demonstrated in Appendix C, LFWA can be successfully estimated with as few as 50 samples with no loss in performance, fully alleviating data access concerns.
Second, our current benchmark evaluates image classification tasks using ResNet-18. In future work, we plan to extend LFWA to vision-language models (such as CLIP) and large language models (LLMs) to verify its parameters' scalability on larger transformers.

\section{Conclusion}
\label{sec:conclusion}
In this paper, we proposed Layer-wise Fisher-Weighted Adaptation (LFWA), a novel, mathematically grounded optimization framework for online test-time model merging. By pre-computing average parameter-level Fisher sensitivities, LFWA dynamically scales layer-wise learning rates to damp updates in sensitive parameter blocks and accelerate them in robust, representational blocks. 
Our empirical results on alternating and block-sequential CIFAR-10 and SVHN streams demonstrate that LFWA consistently outperforms static merging and uniform-rate TTA across both Adam and SGD optimizers, while dramatically accelerating adaptation at low learning rates and reducing sensitivity to global optimization hyperparameters. 
LFWA provides a highly stable, parameter-efficient, and practical optimization foundation for dynamic multi-task systems.

\section*{Impact Statement}
This paper presents work whose goal is to advance the field of Machine Learning. There are many potential societal consequences of our work, none which we feel must be specifically highlighted here.

\bibliography{submission}
\bibliographystyle{icml2026}

%%%%%%%%%%%%%%%%%%%%%%%%%%%%%%%%%%%%%%%%%%%%%%%%%%%%%%%%%%%%%%%%%%%%%%%%%%%%%%%
%%%%%%%%%%%%%%%%%%%%%%%%%%%%%%%%%%%%%%%%%%%%%%%%%%%%%%%%%%%%%%%%%%%%%%%%%%%%%%%
% APPENDIX
%%%%%%%%%%%%%%%%%%%%%%%%%%%%%%%%%%%%%%%%%%%%%%%%%%%%%%%%%%%%%%%%%%%%%%%%%%%%%%%
%%%%%%%%%%%%%%%%%%%%%%%%%%%%%%%%%%%%%%%%%%%%%%%%%%%%%%%%%%%%%%%%%%%%%%%%%%%%%%%
\newpage
\appendix
\onecolumn
\section{Evolution and Distribution of Fisher Information}

In this appendix, we provide additional empirical details regarding the pre-computed layer-wise Fisher sensitivities of the pre-trained ResNet-18 model. Figure~\ref{fig:fisher_dist} illustrates the joint parameter Fisher sensitivities ($\bar{F}_w$) across the ResNet-18 layers.

\begin{figure}[htbp]
\centering
\includegraphics[width=0.7\linewidth]{fisher_dist.pdf}
\caption{Average parameter-level empirical diagonal Fisher Information Matrix (FIM) sensitivities across named parameter layers in ResNet-18. Sensitivities are plotted on a log scale, demonstrating the extreme variance (spanning up to 4 orders of magnitude) between sensitive (early conv/late projection) and robust (intermediate residual) layers.}
\label{fig:fisher_dist}
\end{figure}

As shown in Figure~\ref{fig:fisher_dist}, the classification-adjacent layer (the projection before the classification head) and the initial convolutional blocks exhibit FIM sensitivities that are up to 3 to 4 orders of magnitude larger than intermediate residual blocks. This extreme variance in sensitivity explains why uniform-rate test-time adaptation is inherently unstable: any uniform update step that is large enough to adapt the intermediate blocks inevitably over-updates and damages these highly sensitive structural points, leading to immediate representational degradation. By damping updates in these high-Fisher zones, LFWA preserves structural integrity while unlocking fast and stable learning overall.

\newpage
\section{TTA Merging Coefficient Trajectories}
\label{sec:appendix_trajectories}

To provide deeper insight into the internal optimization dynamics of our framework, we track and visualize the trajectory of the merging coefficients $\lambda_{1, w}$ (corresponding to the CIFAR-10 expert's weight) over the 32 batches of the Alternating stream. We compare standard uniform-rate test-time adaptation ($\alpha=0.0, \eta=0.1$) against our proposed LFWA framework ($\alpha=0.5, \eta=0.001$). Figure~\ref{fig:coefficient_trajectories} displays the results for three representative layers of varying sensitivities:
\begin{enumerate}\setlength{\itemsep}{1pt}\setlength{\parskip}{0pt}\setlength{\parsep}{0pt}
    \item \textbf{Early convolutional layer} (\texttt{conv1.weight}, $F_w \approx 0.0235$), which represents a highly sensitive structural layer.
    \item \textbf{Middle convolutional layer} (\texttt{layer2.0.conv1.weight}, $F_w \approx 0.00036$), representing intermediate sensitivity.
    \item \textbf{Late convolutional layer} (\texttt{layer4.1.conv2.weight}, $F_w \approx 10^{-6}$), representing a robust deep representational layer.
\end{enumerate}

\begin{figure}[htbp]
\centering
\includegraphics[width=0.9\linewidth]{coefficient_trajectories.pdf}
\caption{Evolution of the merging coefficient $\lambda_{1, w}$ for the CIFAR-10 expert over 32 adaptation batches under the Alternating stream. We compare Standard TTA (left) with LFWA (right). While Standard TTA forces highly sensitive early layers to change too rapidly (leading to structural collapse), LFWA adaptively damps sensitive updates while allowing robust layers to adapt with high velocity, stabilizing and accelerating the overall adaptation process.}
\label{fig:coefficient_trajectories}
\end{figure}

As shown in Figure~\ref{fig:coefficient_trajectories} (left), Standard TTA updates all layers using a uniform step size. Consequently, the merging coefficient of the highly sensitive early convolutional layer (\texttt{conv1.weight}, red curve) undergoes massive fluctuations, rapidly changing from its initialization at 0.5 to extreme values. This destabilizes the entire feature extractor and leads to immediate representational degradation and lower overall multi-task performance. 

In contrast, as shown in Figure~\ref{fig:coefficient_trajectories} (right), LFWA scales layer-wise learning rates inversely proportional to their Fisher Information. This adaptively damps the updates of the highly sensitive early convolutional layer, allowing it to adapt slowly and smoothly. Meanwhile, the robust late convolutional layer (\texttt{layer4.1.conv2.weight}, green curve) is assigned a larger effective learning rate, adapting rapidly to find the optimal task blend for each incoming batch. This empirically validates our core hypothesis: LFWA successfully decouples the adaptation speeds of different layers based on their structural sensitivities, preventing representation collapse in crucial feature extraction blocks while accelerating adaptation in representation-rich ones.



\newpage
\section{Additional Ablation Studies and Resource Profiling}
\label{sec:appendix_ablations}

In this appendix, we present the detailed empirical results of our secondary ablation studies and computational resource profiling.

\subsection{Sensitivity to Calibration Set Size}
A potential concern of incorporating pre-computed Fisher Information is the need for calibration data from the experts' training distributions. To evaluate the sensitivity of LFWA to the calibration set size $N_{\text{cal}}$, we estimate the joint layer-wise Fisher sensitivities using varying sample sizes $N_{\text{cal}} \in \{50, 100, 250, 500\}$. We then calculate the Spearman rank correlation and Pearson correlation of these sensitivities against our standard 500-sample FIM. Finally, we evaluate the downstream TTA performance on both Alternating streams (with Adam, using $\eta = 0.001$ and $\alpha = 0.5$) and Block-Sequential streams (with Adam, using $\eta = 0.1$ and $\alpha = 0.2$).

As shown in Table~\ref{tab:calib_ablation}, the joint layer-wise Fisher sensitivities are remarkably robust to the calibration set size. Even with as few as 50 samples, the Spearman rank correlation with the 500-sample FIM is 0.9991, indicating that the rank order of layer-wise sensitivities is extremely stable. Furthermore, the downstream multi-task TTA accuracy remains virtually identical across all sizes (varying by less than 0.5\%). This demonstrates that LFWA is highly data-efficient and can be reliably deployed with minimal (e.g., 50 samples) calibration data, significantly mitigating concerns about training data access.

\begin{table}[htbp]
\caption{Impact of calibration set size $N_{\text{cal}}$ on layer-wise Fisher sensitivities and downstream TTA performance. Correlations are calculated relative to the standard 500-sample FIM.}
\label{tab:calib_ablation}
\vskip 0.15in
\begin{center}
\begin{small}
\resizebox{\columnwidth}{!}{%
\begin{tabular}{ccccc}
\toprule
$N_{\text{cal}}$ & Spearman ($\rho$) & Pearson ($r$) & Alternating (\%) & Sequential (\%) \\
\midrule
50 & 0.9991 & 0.9989 & 79.72 & 81.04 \\
100 & 0.9992 & 0.9992 & 79.81 & 81.24 \\
250 & 0.9997 & 0.9940 & 79.47 & 81.48 \\
500 & 1.0000 & 1.0000 & 79.67 & 81.55 \\
\bottomrule
\end{tabular}%
}
\end{small}
\end{center}
\end{table}

\subsection{Sensitivity to the Number of Adaptation Steps}
In many real-world TTA scenarios, the model can perform multiple optimization steps per incoming batch to more thoroughly adapt the merging coefficients. However, taking multiple steps also increases the risk of representation collapse in sensitive layers. In this section, we study the impact of varying the number of adaptation steps per batch $S \in \{1, 3, 5\}$ using the Adam optimizer across both streams.

As shown in Table~\ref{tab:steps_ablation}, taking more steps generally accelerates adaptation and improves accuracy across all configurations. Most notably, under our LFWA framework, performing multiple steps leads to exceptional performance gains:
First, for LFWA ($\alpha=0.5$, $\eta=0.001$) on the Block-Sequential stream, increasing the steps from 1 to 5 improves accuracy from $81.05\%$ to \textbf{83.51\%} (the highest score achieved on this benchmark). This shows that with a safe preconditioned optimizer, taking more steps unlocks deeper, more effective adaptation without any risk of representation collapse.
Second, under the challenging Alternating stream, LFWA ($\alpha=0.5$, $\eta=0.1$) achieves a major performance jump from $78.12\%$ ($S=1$) to \textbf{81.20\%} ($S=5$). In contrast, standard TTA ($\alpha=0.0$) at the same learning rate plateaus quickly, improving only from $78.82\%$ to $79.40\%$. 
These findings demonstrate that LFWA provides a significantly safer and more scale-receptive optimization landscape, allowing practitioners to leverage multiple local steps to dramatically enhance adaptation accuracy in dynamic environments.

\begin{table*}[htbp]
\caption{Impact of the number of test-time adaptation steps per batch $S \in \{1, 3, 5\}$ on model accuracy (\%) using the Adam optimizer. We compare Standard TTA ($\alpha=0.0$) with LFWA ($\alpha \in \{0.2, 0.5\}$).}
\label{tab:steps_ablation}
\vskip 0.15in
\begin{center}
\begin{small}
\begin{tabular}{lccccccc}
\toprule
& & \multicolumn{3}{c}{\textbf{Alternating Stream (\%)}} & \multicolumn{3}{c}{\textbf{Block-Sequential Stream (\%)}} \\
\cmidrule(r){3-5} \cmidrule(l){6-8}
Method & Learning Rate $\eta$ & $S=1$ & $S=3$ & $S=5$ & $S=1$ & $S=3$ & $S=5$ \\
\midrule
Standard TTA ($\alpha = 0.0$) & 0.001 & 74.73 & 75.51 & 76.51 & 74.82 & 76.23 & 76.96 \\
Standard TTA ($\alpha = 0.0$) & 0.01  & 78.06 & 79.38 & 79.66 & 77.56 & 79.59 & 82.18 \\
Standard TTA ($\alpha = 0.0$) & 0.1   & 78.82 & 79.06 & 79.40 & 81.26 & 81.81 & 82.43 \\
\midrule
LFWA ($\alpha = 0.2$)         & 0.001 & 77.76 & 78.64 & 79.16 & 77.58 & 79.34 & 81.49 \\
LFWA ($\alpha = 0.2$)         & 0.01  & 78.77 & 79.43 & 79.65 & 81.17 & 82.37 & 82.79 \\
LFWA ($\alpha = 0.2$)         & 0.1   & 79.17 & 78.87 & 79.94 & 81.55 & 81.88 & 83.02 \\
\midrule
LFWA ($\alpha = 0.5$)         & 0.001 & 79.67 & 79.43 & 79.63 & 81.05 & 82.17 & \textbf{83.51} \\
LFWA ($\alpha = 0.5$)         & 0.01  & 79.14 & 78.96 & 80.28 & 80.74 & 82.18 & 82.92 \\
LFWA ($\alpha = 0.5$)         & 0.1   & 78.12 & \textbf{80.31} & \textbf{81.20} & 79.83 & 81.27 & 80.75 \\
\bottomrule
\end{tabular}
\end{small}
\end{center}
\end{table*}

\subsection{Sensitivity to the Smoothing Constant $\epsilon$}
The LFWA learning rate scaling formulation scales the global learning rate inversely proportional to $(\bar{F}_l + \epsilon)^\alpha$. The smoothing constant $\epsilon$ serves to ensure numerical stability and prevent division by zero in layers with near-zero average Fisher Information. In this section, we study how the choice of $\epsilon \in \{10^{-10}, 10^{-8}, 10^{-6}, 10^{-4}\}$ impacts the performance of LFWA. We evaluate the Alternating stream (under Adam with $\eta=0.001, \alpha=0.5$) and the Block-Sequential stream (under Adam with $\eta=0.1, \alpha=0.2$).

As shown in Table~\ref{tab:epsilon_ablation}, LFWA is exceptionally robust to the choice of $\epsilon$ over several orders of magnitude. For the Alternating stream, the accuracy remains completely flat at $79.67\%$ for $\epsilon \in \{10^{-10}, 10^{-8}\}$ and only drops slightly to $79.62\%$ at $\epsilon=10^{-6}$. Similarly, on the Sequential stream, accuracy varies by at most $0.15\%$ across $\epsilon \in \{10^{-10}, 10^{-8}, 10^{-6}\}$, showing extreme stability. 
When $\epsilon$ is increased to a relatively large value of $10^{-4}$, it begins to dominate the denominator $(\bar{F}_l + \epsilon)^\alpha$, which mathematically suppresses the differences in preconditioning factors between layers and pulls the optimization back towards a standard uniform-rate update. This causes accuracy to drop slightly to $78.82\%$ on the Alternating stream, which corresponds to the standard TTA accuracy under the same learning rate. This confirms that a reasonably small $\epsilon$ (e.g., $10^{-8}$) is sufficient to enable proper layer-wise preconditioning.

\subsection{Resource Profiling: Memory and Computational Efficiency}
To verify the practical, real-world utility of LFWA, we perform a rigorous resource profiling experiment comparing the test-time GPU memory consumption and per-step adaptation latency of different frameworks. We evaluate on a single NVIDIA H100 GPU using ResNet-18 with a batch size of 64. We compare: (1) \textbf{Static Merging Baseline}: Standard forward inference without any backward pass or parameter optimization. (2) \textbf{Teacher-Free TTA (LFWA / Standard TTA)}: Our framework, where only the single virtual merged student model is loaded into VRAM, and updated using prediction entropy minimization. (3) \textbf{Teacher-Guided TTA}: Emulating expert-guided methods (such as SyMerge \cite{jung2025symerge} or SATA-SBF \cite{submission1}) where the student model and $K=2$ unmerged expert models must be kept in VRAM to compute teacher-guided loss functions.

As shown in Table~\ref{tab:efficiency}, the Static Merging baseline requires 1437.6 MB of GPU memory with an inference latency of 31.0 ms per batch. Our proposed Teacher-Free TTA (LFWA) adds only a minimal memory overhead (+135.7 MB) and has a latency of 97.6 ms, which is highly efficient. 
In contrast, Teacher-Guided TTA requires keeping multiple unmerged expert encoders in memory, which increases GPU memory consumption to 2125.1 MB (a \textbf{35.1\% increase in VRAM} over LFWA) and slows down adaptation to 159.7 ms per batch (a \textbf{63.6\% increase in latency}). 

This quantitative profiling highlights the ``Teacher-Overhead Paradox'' of expert-guided model merging: while keeping multiple expert models in VRAM can help guide adaptation, it fundamentally defeats the primary memory-saving motivation of model merging. By achieving state-of-the-art test-time performance without any expert models in memory, LFWA resolves this paradox, providing a highly scalable and computationally lightweight optimization framework for real-world resource-constrained applications.

\begin{table*}[htbp]
\begin{minipage}{0.48\textwidth}
\centering
\caption{Impact of the smoothing constant $\epsilon$ on the adaptation accuracy (\%) of the LFWA framework using the Adam optimizer.}
\label{tab:epsilon_ablation}
\vskip 0.15in
\begin{small}
\resizebox{\linewidth}{!}{%
\begin{tabular}{ccc}
\toprule
Smoothing Constant $\epsilon$ & Alternating Stream (\%) & Block-Sequential Stream (\%) \\
\midrule
$10^{-10}$ & 79.67 & 81.60 \\
$10^{-8}$ (Default) & 79.67 & 81.55 \\
$10^{-6}$ & 79.62 & 81.70 \\
$10^{-4}$ & 78.82 & 81.39 \\
\bottomrule
\end{tabular}%
}
\end{small}
\end{minipage}
\hfill
\begin{minipage}{0.48\textwidth}
\centering
\caption{GPU Peak Memory (VRAM) and per-step Latency profiling on a single NVIDIA H100 GPU with ResNet-18 (batch size 64). Latency factor represents the slowdown relative to our proposed Teacher-Free TTA (LFWA).}
\label{tab:efficiency}
\vskip 0.15in
\begin{small}
\resizebox{\linewidth}{!}{%
\begin{tabular}{lccc}
\toprule
Method & Peak VRAM (MB) & Latency (ms) & Latency Factor \\
\midrule
Static Merging & 1437.6 & 31.0 & -- \\
Teacher-Free TTA (LFWA) & 1573.3 & 97.6 & 1.0$\times$ \\
Teacher-Guided TTA & 2125.1 & 159.7 & 1.6$\times$ \\
\bottomrule
\end{tabular}%
}
\end{small}
\end{minipage}
\end{table*}

\end{document}
